\documentclass[11pt]{article}

\usepackage[a4paper,margin=1in]{geometry}
\usepackage{amsmath,amssymb,amsthm}
\usepackage{mathtools}
\usepackage{enumitem}
\usepackage{hyperref}

\title{Minimal Dynamic Partitions and Marginal Irreducibility}
\author{}
\date{}

\newcommand{\Obs}{\operatorname{obs}}
\newcommand{\Fin}{\operatorname{Fin}}
\newcommand{\Future}{\operatorname{Future}}
\newcommand{\Sig}{\operatorname{Sig}}
\newcommand{\im}{\operatorname{im}}
\newcommand{\Ker}{\operatorname{Ker}}
\newcommand{\two}{\mathbf{2}}
\newcommand{\id}{\operatorname{id}}

\newtheorem{definition}{Definition}
\newtheorem{proposition}{Proposition}
\newtheorem{remark}{Remark}
\newtheorem{example}{Example}

\begin{document}

\maketitle

\begin{abstract}
We define a relative invariant for a system observed through two marginal interfaces and their joint refinement.
Given a joint fiber \(X=F_{AB}(h)\), two marginal projections
\[
p_A:X\to Y_A,
\qquad
p_B:X\to Y_B,
\]
and a full future-signature map
\[
\Sig:X\to \two^{\Future(h)},
\]
the central object is the simultaneous configuration of three equivalence relations on \(X\):
\[
E_A=\Ker(p_A),
\qquad
E_B=\Ker(p_B),
\qquad
E_\Sig=\{(x,x')\in X\times X:\Sig(x)=\Sig(x')\}.
\]
The invariant is the configuration
\[
\mathcal I_{AB}(h):=[X;E_A,E_B,E_\Sig],
\]
taken up to isomorphism of partition configurations.
Equivalently, it is the intrinsic diagram
\[
X\longrightarrow X/E_A\times X/E_B\times X/E_\Sig.
\]
Its dimension is
\[
\dim \mathcal I_{AB}(h)=|X/E_\Sig|=|\im(\Sig)|.
\]
Marginal irreducibility is the relative-position condition
\[
E_A\nsubseteq E_\Sig,
\qquad
E_B\nsubseteq E_\Sig.
\]
Thus, the invariant is not merely a cardinal. It is a configuration of partitions: the two marginal partitions and the coarsest dynamic quotient preserving the full future signature.
Scientifically, the quotient \(X/E_\Sig\) can be read as the minimal dynamically sufficient state space, while the full invariant \(\mathcal I_{AB}(h)\) records its position relative to the observed marginal state spaces.
\end{abstract}

\section{Interfaces and Observational Equivalence}

Let \(S\) be a state space. An observable interface is a map
\[
\pi:S\to V.
\]

It induces an observational equivalence relation
\[
s\sim_\pi s'
\iff
\pi(s)=\pi(s').
\]

Thus, an interface determines a partition of the state space into observational classes.

\section{Marginal Interfaces and Joint Refinement}

Let
\[
\Obs_A:S\to V_A,
\qquad
\Obs_B:S\to V_B
\]
be two marginal interfaces.

They induce equivalence relations
\[
s\sim_A s'
\iff
\Obs_A(s)=\Obs_A(s'),
\]
and
\[
s\sim_B s'
\iff
\Obs_B(s)=\Obs_B(s').
\]

The joint interface is
\[
\Obs_{AB}:S\to V_A\times V_B,
\]
defined by
\[
\Obs_{AB}(s)=\bigl(\Obs_A(s),\Obs_B(s)\bigr).
\]

It induces
\[
s\sim_{AB}s'
\iff
s\sim_A s'\wedge s\sim_B s'.
\]

Hence
\[
\sim_{AB}=\sim_A\cap\sim_B.
\]

In the refinement order on interfaces,
\[
AB\preceq A,
\qquad
AB\preceq B,
\]
where \(X\preceq Y\) means that \(X\) is finer than \(Y\).

Thus, the joint interface is the meet of the two marginals:
\[
AB=A\wedge B.
\]

\section{Joint Fiber and Marginal Projections}

Fix a prefix or context \(h\).

Let
\[
F_A(h),
\qquad
F_B(h),
\qquad
F_{AB}(h)
\]
be the fibers associated with the interfaces \(A\), \(B\), and \(AB\).

Set
\[
X:=F_{AB}(h),
\qquad
Y_A:=F_A(h),
\qquad
Y_B:=F_B(h).
\]

The joint fiber carries canonical projections
\[
p_A:X\to Y_A,
\qquad
p_B:X\to Y_B.
\]

Define
\[
E_A:=\Ker(p_A),
\qquad
E_B:=\Ker(p_B),
\]
that is
\[
E_A=\{(x,x')\in X\times X\mid p_A(x)=p_A(x')\},
\]
and
\[
E_B=\{(x,x')\in X\times X\mid p_B(x)=p_B(x')\}.
\]

Equivalently,
\[
x\equiv_A x'
\iff
(x,x')\in E_A,
\]
and
\[
x\equiv_B x'
\iff
(x,x')\in E_B.
\]

\section{Future Signatures}

Let
\[
\Future(h)
\]
be the space of admissible future continuations from \(h\).

The full dynamic signature of a point \(x\in X\) is a function
\[
\Sig(x):\Future(h)\to\two.
\]

Equivalently,
\[
\Sig:X\to \two^{\Future(h)}.
\]

For each \(u\in\Future(h)\), the coordinate
\[
\Sig(x)(u)\in\two
\]
records the dynamic compatibility of \(x\) with the future continuation \(u\).

Thus, the future signature is the complete dynamic profile
\[
\Sig(x)=\bigl(\Sig(x)(u)\bigr)_{u\in\Future(h)}.
\]

\section{The Coarsest Dynamic Quotient}

The signature map induces an equivalence relation
\[
x\sim_\Sig x'
\iff
\Sig(x)=\Sig(x').
\]

Let
\[
Q_\Sig:=X/{\sim_\Sig}.
\]

Let
\[
q_\Sig:X\twoheadrightarrow Q_\Sig
\]
be the canonical quotient map.

Then
\[
q_\Sig(x)=q_\Sig(x')
\iff
\Sig(x)=\Sig(x').
\]

Define
\[
E_\Sig:=\Ker(q_\Sig).
\]

Thus,
\[
E_\Sig=\{(x,x')\in X\times X\mid \Sig(x)=\Sig(x')\}.
\]

The relation \(E_\Sig\) is the coarsest dynamic quotient of the joint fiber preserving the full future signature.

\section{The Relative Invariant}

\begin{definition}[Isomorphism of configurations]
Two configurations
\[
[X;E_A,E_B,E_\Sig]
\qquad\text{and}\qquad
[X';E_A',E_B',E_\Sig']
\]
are isomorphic if there exists a bijection
\[
\varphi:X\to X'
\]
such that, for all \(x,x'\in X\),
\[
(x,x')\in E_A
\iff
(\varphi(x),\varphi(x'))\in E_A',
\]
\[
(x,x')\in E_B
\iff
(\varphi(x),\varphi(x'))\in E_B',
\]
and
\[
(x,x')\in E_\Sig
\iff
(\varphi(x),\varphi(x'))\in E_\Sig'.
\]
\end{definition}

\begin{definition}[Dynamic-interface invariant]
The dynamic-interface invariant at \(h\) is the isomorphism class
\[
\boxed{
\mathcal I_{AB}(h):=[X;E_A,E_B,E_\Sig].
}
\]

Equivalently, it is the intrinsic diagram
\[
\boxed{
X\longrightarrow X/E_A\times X/E_B\times X/E_\Sig
}
\]
up to isomorphism.
\end{definition}

This invariant records the relative position of three partitions on the same set \(X\):
\[
E_A=\Ker(p_A),
\qquad
E_B=\Ker(p_B),
\qquad
E_\Sig=\Ker(q_\Sig).
\]

The quotient \(X/E_\Sig\) gives the dynamic partition. The full invariant keeps track of how this dynamic partition sits relative to the two marginal partitions.

\section{Dimension}

Assume that
\[
|\im(\Sig)|<\infty.
\]

Then
\[
Q_\Sig\cong \im(\Sig),
\]
and hence
\[
|Q_\Sig|=|\im(\Sig)|.
\]

\begin{definition}[Dimension]
The dimension of the invariant is
\[
\boxed{
\dim\mathcal I_{AB}(h):=|X/E_\Sig|=|Q_\Sig|=|\im(\Sig)|.
}
\]
\end{definition}

Thus, the dimension is the number of distinct future signatures realized on the joint fiber.

The full invariant is the configuration
\[
[X;E_A,E_B,E_\Sig],
\]
while the dimension is only its cardinal shadow.

\section{Finite Mediators}

A finite mediator for the signature is a map
\[
M:X\to \Fin(m)
\]
such that \(\Sig\) factors through \(M\):
\[
\Sig=\widehat{\Sig}_M\circ M
\]
for some map
\[
\widehat{\Sig}_M:\Fin(m)\to \two^{\Future(h)}.
\]

Thus, \(M\) is a finite coding of points in \(X\) sufficient to reconstruct their full future signatures.

\begin{proposition}[Minimal dynamic quotient]
Let \(X\) be a set and let
\[
\Sig:X\to Z
\]
be a map with finite image.
Let
\[
q_\Sig:X\twoheadrightarrow Q_\Sig=X/{\sim_\Sig}
\]
where
\[
x\sim_\Sig x'
\iff
\Sig(x)=\Sig(x').
\]
Then \(q_\Sig\) is the coarsest quotient of \(X\) through which \(\Sig\) factors.

Moreover, every finite mediator
\[
M:X\to\Fin(m)
\]
through which \(\Sig\) factors satisfies
\[
m\ge |Q_\Sig|.
\]

If
\[
m=|Q_\Sig|,
\]
then \(M\) is equivalent to \(q_\Sig\) up to a bijection of labels.
\end{proposition}

\begin{proof}
Since \(q_\Sig(x)=q_\Sig(x')\) exactly when \(\Sig(x)=\Sig(x')\), the map \(\Sig\) factors through \(q_\Sig\):
\[
\Sig=\widehat{\Sig}_q\circ q_\Sig
\]
where
\[
\widehat{\Sig}_q([x])=\Sig(x).
\]

Let
\[
M:X\to\Fin(m)
\]
be any finite mediator such that
\[
\Sig=\widehat{\Sig}_M\circ M.
\]

If
\[
M(x)=M(x'),
\]
then
\[
\Sig(x)=\widehat{\Sig}_M(M(x))
=
\widehat{\Sig}_M(M(x'))
=
\Sig(x').
\]

Hence every fiber of \(M\) is contained in a fiber of \(\Sig\). Therefore \(M\) must distinguish all distinct elements of \(\im(\Sig)\), so
\[
|\im(M)|\ge |\im(\Sig)|=|Q_\Sig|.
\]

Since
\[
|\im(M)|\le m,
\]
we get
\[
m\ge |Q_\Sig|.
\]

If equality holds, then the fibers of \(M\) coincide exactly with the fibers of \(\Sig\). Therefore \(M\) differs from \(q_\Sig\) only by a bijective relabeling of the quotient classes.
\end{proof}

\section{Descent Along a Marginal Projection}

Let
\[
C\in\{A,B\}.
\]

Write
\[
p_C:X\to Y_C
\]
for the corresponding marginal projection and
\[
E_C:=\Ker(p_C).
\]

\begin{definition}[Descent]
The dynamic quotient \(q_\Sig\) descends to the marginal \(C\) when there exists a map
\[
\rho_C:Y_C\to Q_\Sig
\]
such that
\[
q_\Sig=\rho_C\circ p_C.
\]

We write
\[
q_\Sig\leq C
\]
when this descent exists.
\end{definition}

\begin{proposition}[Marginal descent criterion]
For \(C\in\{A,B\}\), the following are equivalent:
\[
q_\Sig\leq C,
\]
\[
\Ker(p_C)\subseteq \Ker(q_\Sig),
\]
and
\[
\Sig \text{ factors through } p_C.
\]
\end{proposition}

\begin{proof}
The condition
\[
q_\Sig\leq C
\]
means that there exists
\[
\rho_C:Y_C\to Q_\Sig
\]
such that
\[
q_\Sig=\rho_C\circ p_C.
\]

This is equivalent to constancy of \(q_\Sig\) on every fiber of \(p_C\), namely
\[
p_C(x)=p_C(x')
\Rightarrow
q_\Sig(x)=q_\Sig(x').
\]

In kernel notation, this is exactly
\[
\Ker(p_C)\subseteq\Ker(q_\Sig).
\]

Since
\[
q_\Sig(x)=q_\Sig(x')
\iff
\Sig(x)=\Sig(x'),
\]
the condition becomes
\[
p_C(x)=p_C(x')
\Rightarrow
\Sig(x)=\Sig(x').
\]

This is exactly the condition for the existence of a map
\[
\widehat{\Sig}_C:Y_C\to \two^{\Future(h)}
\]
such that
\[
\Sig=\widehat{\Sig}_C\circ p_C.
\]
\end{proof}

\section{Marginal Irreducibility}

\begin{definition}[Marginal irreducibility]
The invariant
\[
\mathcal I_{AB}(h)=[X;E_A,E_B,E_\Sig]
\]
is marginally irreducible when
\[
\boxed{
E_A\nsubseteq E_\Sig
\qquad
\text{and}
\qquad
E_B\nsubseteq E_\Sig.
}
\]

Equivalently,
\[
\boxed{
\Ker(p_A)\nsubseteq\Ker(q_\Sig)
\qquad
\text{and}
\qquad
\Ker(p_B)\nsubseteq\Ker(q_\Sig).
}
\]
\end{definition}

This says that each marginal identifies at least one pair of joint-fiber points that the future signature separates.

Explicitly, irreducibility along \(A\) means:
\[
\exists x,x'\in X,
\qquad
p_A(x)=p_A(x'),
\qquad
q_\Sig(x)\neq q_\Sig(x').
\]

Since \(q_\Sig\) classifies future signatures, this is equivalent to:
\[
\exists x,x'\in X,
\qquad
p_A(x)=p_A(x'),
\qquad
\Sig(x)\neq\Sig(x').
\]

Equivalently:
\[
\exists x,x'\in X,
\quad
\exists u\in\Future(h),
\qquad
p_A(x)=p_A(x'),
\qquad
\Sig(x)(u)\neq\Sig(x')(u).
\]

The same formulation applies to \(B\).

\section{Diagonal Obstruction}

A diagonal obstruction along \(A\) is the existence of
\[
x,x'\in X
\]
such that
\[
p_A(x)=p_A(x')
\]
and
\[
\Sig(x)\neq\Sig(x').
\]

Equivalently,
\[
E_A\nsubseteq E_\Sig.
\]

Thus,
\[
\boxed{
\text{diagonal obstruction along }A
\iff
q_\Sig \text{ descends to } A \text{ fails}
}
\]
or, in kernel form,
\[
\boxed{
\text{diagonal obstruction along }A
\iff
\Ker(p_A)\nsubseteq\Ker(q_\Sig).
}
\]

Likewise,
\[
\boxed{
\text{diagonal obstruction along }B
\iff
\Ker(p_B)\nsubseteq\Ker(q_\Sig).
}
\]

Hence diagonal obstruction is the witness form of marginal irreducibility.

\section{Relation to Boolean Step Tests}

For a single future
\[
u\in\Future(h),
\]
define the Boolean step test
\[
T_u:X\to\two,
\qquad
T_u(x)=\Sig(x)(u).
\]

A step-level obstruction along \(A\) is:
\[
p_A(x)=p_A(x'),
\qquad
T_u(x)\neq T_u(x').
\]

This shows that the single Boolean test \(T_u\) separates a pair collapsed by \(p_A\).

The full future signature collects the entire family of tests:
\[
\Sig(x)=\bigl(T_u(x)\bigr)_{u\in\Future(h)}.
\]

The dynamic quotient
\[
q_\Sig:X\twoheadrightarrow Q_\Sig
\]
is therefore the common coarsest dynamic quotient induced by all step tests at once.

For a single Boolean predicate
\[
T:X\to\two,
\]
the quotient has at most two classes. For the full signature
\[
\Sig:X\to\two^{\Future(h)},
\]
the quotient size is
\[
|Q_\Sig|=|\im(\Sig)|.
\]

This cardinal can exceed \(2\), even though each coordinate \(T_u\) is Boolean.

\section{A Finite Example}

Let
\[
X=\{00,01,10,11\}.
\]

Define binary marginal projections
\[
p_A(ab)=a,
\qquad
p_B(ab)=b.
\]

Then the marginal equivalence relations are
\[
E_A=\bigl\{\{00,01\},\{10,11\}\bigr\},
\]
and
\[
E_B=\bigl\{\{00,10\},\{01,11\}\bigr\}.
\]

Let
\[
\Future(h)=\{u_1,u_2\},
\]
and define the future signature by
\[
\Sig(00)=(0,0),
\qquad
\Sig(01)=(0,1),
\]
\[
\Sig(10)=(1,0),
\qquad
\Sig(11)=(1,1).
\]

Then all four future signatures are distinct. Hence
\[
E_\Sig=\Delta_X,
\]
where \(\Delta_X\) is the equality relation on \(X\).

Thus,
\[
Q_\Sig\cong X,
\qquad
\dim\mathcal I_{AB}(h)=4.
\]

However,
\[
E_A\nsubseteq E_\Sig,
\qquad
E_B\nsubseteq E_\Sig.
\]

Indeed, \(00\) and \(01\) are identified by \(p_A\), but
\[
\Sig(00)\neq\Sig(01).
\]

Similarly, \(00\) and \(10\) are identified by \(p_B\), but
\[
\Sig(00)\neq\Sig(10).
\]

Therefore the invariant
\[
\mathcal I_{AB}(h)=[X;E_A,E_B,E_\Sig]
\]
is marginally irreducible, with
\[
\dim\mathcal I_{AB}(h)=4.
\]

This example shows that two binary marginals can induce a joint dynamic-signature partition with four classes, and that the irreducibility is a property of the relative position of the partitions, not merely of the cardinality.

\section{Main Statement}

The invariant associated with the two marginal interfaces and the future signature is:
\[
\boxed{
\mathcal I_{AB}(h):=[X;E_A,E_B,E_\Sig].
}
\]

Equivalently:
\[
\boxed{
\mathcal I_{AB}(h)
=
\left[
X\longrightarrow X/E_A\times X/E_B\times X/E_\Sig
\right].
}
\]

Its dimension is:
\[
\boxed{
\dim\mathcal I_{AB}(h)=|X/E_\Sig|=|\im(\Sig)|.
}
\]

It is marginally irreducible exactly when:
\[
\boxed{
E_A\nsubseteq E_\Sig,
\qquad
E_B\nsubseteq E_\Sig.
}
\]

In kernel notation:
\[
\boxed{
\Ker(p_A)\nsubseteq\Ker(q_\Sig),
\qquad
\Ker(p_B)\nsubseteq\Ker(q_\Sig).
}
\]

In words:

\begin{quote}
The invariant is the relative configuration of the two marginal partitions and the coarsest dynamic-signature partition on the same joint fiber.
\end{quote}

\begin{quote}
Marginal irreducibility means that each marginal projection collapses at least one pair that the full future signature separates.
\end{quote}

\section{Conceptual Chain}

The conceptual chain is:
\[
\text{marginal projection}
\Rightarrow
\text{identification}
\Rightarrow
\text{future-signature separation}
\Rightarrow
\text{diagonal obstruction}
\Rightarrow
\text{coarsest dynamic quotient}
\Rightarrow
\text{relative partition invariant}
\Rightarrow
\text{marginal irreducibility}.
\]

Equivalently:

\begin{quote}
A marginal projection collapses points.  
The future signature separates them.  
The quotient by future-signature equality gives the coarsest dynamic quotient.  
The invariant is the configuration of this dynamic partition relative to the marginal partitions.  
Marginal irreducibility is the failure of inclusion
\[
E_A\subseteq E_\Sig
\qquad
\text{and}
\qquad
E_B\subseteq E_\Sig.
\]
\end{quote}

\section{Scientific Interpretation: Minimal Dynamic State Variables}

The invariant
\[
\mathcal I_{AB}(h)=[X;E_A,E_B,E_\Sig]
\]
can be interpreted as a formal device for identifying the minimal state variable required to preserve the future dynamics.

The marginal partitions
\[
E_A=\Ker(p_A),
\qquad
E_B=\Ker(p_B)
\]
represent the information retained by the two partial observations. The dynamic partition
\[
E_\Sig
\]
represents the distinctions that are required to preserve the full future-signature profile.

Thus the quotient
\[
X/E_\Sig
\]
is the minimal dynamic state space associated with the future signature. It groups together exactly those present states that have identical future behavior (as measured by \(\Sig\)).

The relative invariant \(\mathcal I_{AB}(h)\) records how this dynamically sufficient state space sits relative to the observed marginal state spaces \(X/E_A\) and \(X/E_B\).

In scientific terms, this construction separates three notions:
\[
\text{available observation},
\qquad
\text{marginal information},
\qquad
\text{dynamically sufficient state}.
\]

Marginal irreducibility,
\[
E_A\nsubseteq E_\Sig,
\qquad
E_B\nsubseteq E_\Sig,
\]
means that each marginal observation collapses distinctions that are required for dynamic prediction.

Consequently, the invariant identifies the minimal dynamic partition that closes prediction at the joint level, together with the precise obstruction to recovering that partition from either marginal observation alone.

This interpretation is relevant whenever one seeks a coarse-grained state variable that preserves future behavior: statistical physics, dynamical systems, biological state spaces, neuroscience, epidemiology, control, and representation learning.

\section{Status of the Construction}

The quotient construction itself is standard: every map induces a quotient by equality of its values.

The object defined here is the relative configuration:
\[
[X;E_A,E_B,E_\Sig],
\]
where
\[
E_A=\Ker(p_A),
\qquad
E_B=\Ker(p_B),
\qquad
E_\Sig=\Ker(q_\Sig).
\]

Thus, the invariant is a configuration of partitions on the joint fiber, while its dimension is the cardinal
\[
|X/E_\Sig|=|\im(\Sig)|.
\]

The irreducibility conditions
\[
E_A\nsubseteq E_\Sig,
\qquad
E_B\nsubseteq E_\Sig
\]
express the precise sense in which the coarsest dynamic partition is joint relative to the two marginal projections.

\end{document}
